% -------------------------------------------------
% فصل 2
% -------------------------------------------------

\projectchapter{مبانی نظری و مفاهیم پایه}{
در این فصل، مفاهیم نظری مورد نیاز برای درک روش‌های استفاده‌شده در پروژه بررسی می‌شوند. ابتدا مسئله شناسایی زبان از متن و گفتار و دو رویکرد یادگیری نظارت‌شده و بدون نظارت معرفی می‌شوند. سپس مبانی استخراج ویژگی از متن و گفتار، الگوریتم‌های طبقه‌بندی و خوشه‌بندی، روش‌های آماده‌سازی و کاهش ابعاد ویژگی‌ها، اصول اعتبارسنجی و معیارهای ارزیابی تشریح خواهند شد. جزئیات مجموعه‌داده‌ها و پیش‌پردازش در فصل سوم و تنظیمات اجرایی، ابرپارامترها و نتایج آزمایش‌ها در فصل چهارم ارائه می‌شوند.
}


% =================================================
\section{مقدمه}
% =================================================

شناسایی خودکار زبان یکی از مسائل مهم در حوزه پردازش زبان طبیعی و پردازش گفتار است. هدف این مسئله، تعیین زبان یک نمونه ورودی براساس الگوهای قابل استخراج از آن است. ورودی سامانه می‌تواند یک متن یا یک سیگنال گفتاری باشد.

اگر مجموعه زبان‌های مورد بررسی به‌صورت

\begin{equation}
L=\{L_1,L_2,\ldots,L_K\}
\end{equation}

تعریف شود، هدف یک سامانه شناسایی زبان، یادگیری تابعی است که برای نمونه ورودی $x$ یکی از زبان‌های مجموعه $L$ را انتخاب کند:

\begin{equation}
f(x)\longrightarrow L_i
\end{equation}

اگرچه هدف نهایی در مسیرهای متن و گفتار یکسان است، ماهیت ورودی و روش نمایش عددی آن‌ها متفاوت است. در مسیر متن، الگوهای کاراکتری و واژگانی اهمیت دارند؛ درحالی‌که در مسیر گفتار، اطلاعات آوایی، زمانی، طیفی و انرژی سیگنال مبنای مدل‌سازی قرار می‌گیرند.


% =================================================
\section{شناسایی زبان از متن و گفتار}
% =================================================

\subsection{شناسایی زبان از متن}

در شناسایی زبان مبتنی بر متن\footnote{\lr{Text Language Identification}}، هدف تعیین زبان یک متن براساس ساختار نوشتاری آن است. زبان‌ها می‌توانند از نظر نوع خط، ترکیب کاراکترها، ساختار واژه‌ها، فراوانی واژگان و الگوهای محلی متن با یکدیگر تفاوت داشته باشند.

حتی زبان‌هایی که از الفبای مشابه استفاده می‌کنند، معمولاً الگوهای متفاوتی در کنار هم قرارگرفتن حروف، پسوندها، پیشوندها و واژگان دارند. ازاین‌رو، نمایش‌های آماری مبتنی بر کاراکتر و کلمه می‌توانند اطلاعات مناسبی برای شناسایی زبان فراهم کنند.


\subsection{شناسایی زبان از گفتار}

در شناسایی زبان از گفتار\footnote{\lr{Spoken Language Identification}}، ورودی سامانه یک سیگنال صوتی است. یک سیگنال صوتی دیجیتال را می‌توان به‌صورت دنباله‌ای از نمونه‌ها نمایش داد:

\begin{equation}
x[n], \qquad n=0,1,\ldots,N-1
\end{equation}

که در آن $x[n]$ دامنه سیگنال در نمونه $n$ است.

استفاده مستقیم از تمام نمونه‌های سیگنال برای مدل‌های کلاسیک یادگیری ماشین مناسب نیست؛ زیرا فایل‌ها می‌توانند طول‌های متفاوت داشته باشند و هر فایل شامل تعداد زیادی نمونه باشد. بنابراین، هر فایل گفتاری باید به یک بردار ویژگی با طول ثابت تبدیل شود. این بردار می‌تواند ویژگی‌های طیفی، زمانی و انرژی سیگنال را خلاصه کند.


% =================================================
\section{یادگیری ماشین در مسئله شناسایی زبان}
% =================================================

در این پروژه از دو رویکرد اصلی یادگیری ماشین استفاده شده است:

\begin{itemize}

	\item یادگیری نظارت‌شده؛

	\item یادگیری بدون نظارت.

\end{itemize}


% -------------------------------------------------
\subsection{یادگیری نظارت‌شده}
% -------------------------------------------------

در یادگیری نظارت‌شده\footnote{\lr{Supervised Learning}}، برای هر نمونه ورودی یک برچسب صحیح وجود دارد. مجموعه داده آموزشی را می‌توان به‌صورت زیر نمایش داد:

\begin{equation}
D=
\{(x_1,y_1),(x_2,y_2),\ldots,(x_N,y_N)\}
\end{equation}

که در آن:

\begin{itemize}

	\item $x_i$ بردار ویژگی نمونه $i$؛

	\item $y_i$ برچسب زبان نمونه $i$؛

	\item $N$ تعداد نمونه‌ها است.

\end{itemize}

مدل با استفاده از نمونه‌های آموزشی، رابطه میان ویژگی‌ها و زبان واقعی را یاد می‌گیرد و سپس زبان نمونه‌های جدید را پیش‌بینی می‌کند. مسئله طبقه‌بندی\footnote{\lr{Classification}} در این گروه قرار می‌گیرد.


% -------------------------------------------------
\subsection{یادگیری بدون نظارت}
% -------------------------------------------------

در یادگیری بدون نظارت\footnote{\lr{Unsupervised Learning}}، برچسب واقعی نمونه‌ها هنگام برازش الگوریتم در اختیار آن قرار نمی‌گیرد. در این حالت، مجموعه داده به‌صورت

\begin{equation}
D=\{x_1,x_2,\ldots,x_N\}
\end{equation}

تعریف می‌شود.

الگوریتم تلاش می‌کند نمونه‌های مشابه را براساس ساختار ویژگی‌هایشان در گروه‌هایی قرار دهد. خوشه‌بندی\footnote{\lr{Clustering}} نمونه‌ای از یادگیری بدون نظارت است. در مسئله شناسایی زبان، خوشه‌بندی می‌تواند نشان دهد که آیا داده‌ها بدون استفاده از برچسب زبان دارای گروه‌بندی طبیعی هستند یا خیر.


% =================================================
\section{استخراج ویژگی از داده‌های متنی}
% =================================================

الگوریتم‌های کلاسیک یادگیری ماشین نمی‌توانند متن خام را مستقیماً پردازش کنند؛ بنابراین، متن باید به یک نمایش عددی تبدیل شود. در این پروژه، این نمایش با استفاده از الگوهای کاراکتری و واژگانی و وزن‌دهی
\lr{TF--IDF}
ساخته می‌شود.


% -------------------------------------------------
\subsection{مفهوم \lr{N-Gram}}
% -------------------------------------------------

\lr{N-Gram}
به دنباله‌ای شامل $n$ عنصر متوالی گفته می‌شود. این عناصر می‌توانند کاراکتر یا کلمه باشند.

برای مثال، برخی از
\lr{Character Bigram}
های واژه

\begin{center}
\lr{language}
\end{center}

عبارت‌اند از:

\begin{center}
\lr{la, an, ng, gu, ua, ag, ge}
\end{center}

همچنین، برای عبارت

\begin{center}
\lr{this is a language}
\end{center}

برخی
\lr{Word Bigram}
ها عبارت‌اند از:

\begin{center}
\lr{this is, is a, a language}
\end{center}

ویژگی‌های کاراکتری می‌توانند الگوهای مربوط به شکل واژه‌ها، ترکیب حروف و نوع خط را ثبت کنند. ویژگی‌های واژگانی نیز اطلاعات مربوط به واژه‌ها و ترتیب محلی آن‌ها را در اختیار مدل قرار می‌دهند. استفاده هم‌زمان از این دو نوع ویژگی، نمایش کامل‌تری از ساختار متن ایجاد می‌کند.


% -------------------------------------------------
\subsection{روش \lr{TF--IDF}}
% -------------------------------------------------

روش
\lr{TF--IDF}\footnote{\lr{Term Frequency--Inverse Document Frequency}}
یکی از روش‌های متداول برای تبدیل متن به نمایش عددی است. این روش از دو مؤلفه اصلی تشکیل می‌شود.


\subsubsection{\lr{Term Frequency}}

مؤلفه
\lr{Term Frequency}
میزان تکرار ویژگی $t$ در سند $d$ را اندازه‌گیری می‌کند. در حالت ساده می‌توان نوشت:

\begin{equation}
TF(t,d)=\mathrm{count}(t,d)
\end{equation}

در حالت لگاریتمی یا
\lr{Sublinear TF}،
برای فراوانی‌های غیرصفر داریم:

\begin{equation}
TF(t,d)=1+\log(\mathrm{count}(t,d))
\end{equation}

استفاده از لگاریتم مانع از آن می‌شود که تکرار بسیار زیاد یک ویژگی اثر نامتناسبی بر نمایش سند داشته باشد.


\subsubsection{\lr{Inverse Document Frequency}}

مؤلفه
\lr{IDF}
نشان می‌دهد یک ویژگی در کل مجموعه اسناد تا چه اندازه خاص یا عمومی است. فرم هموارشده آن را می‌توان به‌صورت زیر نمایش داد:

\begin{equation}
IDF(t)
=
\log
\left(
\frac{1+N}{1+df(t)}
\right)
+1
\end{equation}

که در آن:

\begin{itemize}

	\item $N$ تعداد کل اسناد؛

	\item $df(t)$ تعداد اسنادی است که ویژگی $t$ در آن‌ها وجود دارد.

\end{itemize}

در نهایت، وزن
\lr{TF--IDF}
از رابطه زیر به دست می‌آید:

\begin{equation}
TFIDF(t,d)
=
TF(t,d)\times IDF(t)
\end{equation}

در نتیجه، ویژگی‌هایی که در یک سند پرتکرار ولی در کل مجموعه نسبتاً کمیاب هستند، وزن بیشتری دریافت می‌کنند.


% -------------------------------------------------
\subsection{ترکیب ویژگی‌های کاراکتری و واژگانی}
% -------------------------------------------------

اگر ماتریس ویژگی‌های کاراکتری با $X_c$ و ماتریس ویژگی‌های واژگانی با $X_w$ نمایش داده شود، نمایش ترکیبی متن را می‌توان به‌صورت الحاق افقی دو ماتریس تعریف کرد:

\begin{equation}
X=
\left[
X_c
\quad
X_w
\right]
\end{equation}

این ترکیب موجب می‌شود مدل هم از اطلاعات سطح کاراکتر و هم از اطلاعات سطح کلمه استفاده کند. تنظیمات دقیق بردارسازها و ابعاد نمایش نهایی در فصل سوم ارائه خواهند شد.


% -------------------------------------------------
\subsection{انتخاب ویژگی با \lr{Chi-Square}}
% -------------------------------------------------

در مسائل طبقه‌بندی متن، فضای اولیه ویژگی‌ها می‌تواند بسیار بزرگ باشد. یکی از روش‌های انتخاب ویژگی، آزمون
\lr{Chi-Square}
است که میزان وابستگی آماری هر ویژگی به برچسب کلاس را بررسی می‌کند.

رابطه کلی آماره کای‌دو به‌صورت زیر است:

\begin{equation}
\chi^2
=
\sum
\frac{(O-E)^2}{E}
\end{equation}

که در آن $O$ مقدار مشاهده‌شده و $E$ مقدار مورد انتظار است. مقدار بزرگ‌تر
$\chi^2$
نشان‌دهنده وابستگی بیشتر ویژگی به کلاس است.

ازآنجاکه این روش به برچسب زبان نیاز دارد، فقط در مسیر نظارت‌شده طبقه‌بندی متن قابل استفاده است. استفاده از آن در ساخت نمایش بدون نظارت خوشه‌بندی موجب ورود اطلاعات برچسب و نقض ماهیت خوشه‌بندی می‌شود. تعداد دقیق ویژگی‌های انتخاب‌شده و نحوه برازش انتخاب‌گر در فصل سوم تشریح خواهند شد.


% =================================================
\section{استخراج ویژگی از داده‌های گفتاری}
% =================================================

هدف استخراج ویژگی گفتاری، تبدیل سیگنال خام با طول متغیر به یک بردار عددی با طول ثابت است. برای این منظور، ویژگی‌های فریم‌محور محاسبه و سپس با آماره‌هایی مانند میانگین و انحراف معیار در سطح کل فایل خلاصه می‌شوند. ترکیب و تعداد دقیق ویژگی‌های نسخه نهایی در فصل سوم ارائه خواهند شد.


% -------------------------------------------------
\subsection{ضرایب \lr{MFCC}}
% -------------------------------------------------

ضرایب
\lr{MFCC}\footnote{\lr{Mel-Frequency Cepstral Coefficients}}
از پرکاربردترین ویژگی‌ها در پردازش گفتار هستند. مقیاس
\lr{Mel}
برای نزدیک‌کردن نمایش فرکانس به نحوه ادراک شنوایی انسان استفاده می‌شود. تبدیل تقریبی فرکانس معمولی به مقیاس
\lr{Mel}
به‌صورت زیر است:

\begin{equation}
Mel(f)
=
2595
\log_{10}
\left(
1+\frac{f}{700}
\right)
\end{equation}

پس از محاسبه طیف، انرژی در بانک فیلترهای
\lr{Mel}
محاسبه می‌شود. سپس لگاریتم انرژی‌ها و تبدیل کسینوسی گسسته برای تولید ضرایب به کار می‌روند. یک فرم کلی برای ضریب شماره $n$ عبارت است از:

\begin{equation}
c_n
=
\sum_{k=1}^{M}
\log(M_k)
\cos
\left[
\frac{\pi n}{M}
\left(
k-\frac{1}{2}
\right)
\right]
\end{equation}

که در آن $M_k$ انرژی فیلتر شماره $k$ و $M$ تعداد فیلترها است.

چون ضرایب
\lr{MFCC}
برای فریم‌های متعدد محاسبه می‌شوند، می‌توان مقدار میانگین و انحراف معیار هر ضریب را برای ساخت یک نمایش با طول ثابت محاسبه کرد:

\begin{equation}
\mu_i
=
\frac{1}{T}
\sum_{t=1}^{T}
c_i(t)
\end{equation}

\begin{equation}
\sigma_i
=
\sqrt{
\frac{1}{T}
\sum_{t=1}^{T}
(c_i(t)-\mu_i)^2
}
\end{equation}


% -------------------------------------------------
\subsection{\lr{Delta MFCC} و \lr{Delta-Delta MFCC}}
% -------------------------------------------------

ضرایب
\lr{MFCC}
وضعیت طیفی هر فریم را توصیف می‌کنند، اما تغییرات این ضرایب در طول زمان نیز حاوی اطلاعات گفتاری است. ویژگی
\lr{Delta MFCC}
سرعت تغییر ضرایب و
\lr{Delta-Delta MFCC}
تغییرات مرتبه دوم آن‌ها را نمایش می‌دهد.

یک تقریب متداول برای مشتق مرتبه اول به‌صورت زیر است:

\begin{equation}
\Delta c_t
=
\frac{
\sum_{n=1}^{R}
n(c_{t+n}-c_{t-n})
}{
2\sum_{n=1}^{R}n^2
}
\end{equation}

که در آن $R$ عرض پنجره زمانی است. مشتق مرتبه دوم با اعمال دوباره همین مفهوم بر ضرایب مشتق مرتبه اول به دست می‌آید. این ویژگی‌ها اطلاعات مربوط به پویایی زمانی گفتار را تکمیل می‌کنند.


% -------------------------------------------------
\subsection{\lr{Zero Crossing Rate}}
% -------------------------------------------------

نرخ عبور از صفر یا
\lr{Zero Crossing Rate}
میزان تغییر علامت سیگنال را در یک بازه زمانی اندازه‌گیری می‌کند. برای یک فریم شامل $N$ نمونه می‌توان نوشت:

\begin{equation}
ZCR
=
\frac{1}{2(N-1)}
\sum_{n=1}^{N-1}
\left|
\mathrm{sgn}(x[n])
-
\mathrm{sgn}(x[n-1])
\right|
\end{equation}

این ویژگی می‌تواند اطلاعاتی درباره میزان تغییرات سریع و ماهیت نویزی یا آوایی سیگنال فراهم کند.


% -------------------------------------------------
\subsection{ویژگی‌های طیفی}
% -------------------------------------------------

ویژگی‌های طیفی برای توصیف شکل و توزیع انرژی در حوزه فرکانس استفاده می‌شوند.


\subsubsection{\lr{Spectral Centroid}}

مرکز طیفی، مرکز جرم طیف فرکانسی را نشان می‌دهد:

\begin{equation}
SC
=
\frac{
\sum_k f_k S_k
}{
\sum_k S_k
}
\end{equation}

که در آن $f_k$ فرکانس و $S_k$ مقدار طیفی متناظر است. مقدار بالاتر این ویژگی به تمرکز بیشتر انرژی در فرکانس‌های بالا اشاره دارد.


\subsubsection{\lr{Spectral Rolloff}}

افت طیفی، فرکانسی است که نسبت مشخصی از انرژی طیف در زیر آن قرار می‌گیرد. اگر $f_r$ فرکانس
\lr{Rolloff}
باشد، داریم:

\begin{equation}
\sum_{f_k\leq f_r}S_k
\geq
\alpha
\sum_k S_k
\end{equation}

که در آن $\alpha$ نسبت انرژی مورد نظر است.


\subsubsection{\lr{Spectral Bandwidth}}

پهنای باند طیفی میزان پراکندگی طیف در اطراف مرکز طیفی را توصیف می‌کند:

\begin{equation}
BW
=
\sqrt{
\frac{
\sum_k S_k(f_k-SC)^2
}{
\sum_k S_k
}
}
\end{equation}


\subsubsection{\lr{Spectral Flatness}}

تختی طیفی نسبت میانگین هندسی به میانگین حسابی مقادیر طیف است:

\begin{equation}
SF
=
\frac{
\left(\prod_{k=1}^{M}S_k\right)^{1/M}
}{
\frac{1}{M}\sum_{k=1}^{M}S_k
}
\end{equation}

این ویژگی می‌تواند میزان تخت، نویزگونه یا قله‌داربودن طیف را توصیف کند.


% -------------------------------------------------
\subsection{\lr{RMS Energy}}
% -------------------------------------------------

انرژی
\lr{RMS}\footnote{\lr{Root Mean Square}}
معیاری برای توصیف سطح انرژی سیگنال است:

\begin{equation}
RMS
=
\sqrt{
\frac{1}{N}
\sum_{n=0}^{N-1}x[n]^2
}
\end{equation}


% -------------------------------------------------
\subsection{مدت‌زمان فایل صوتی}
% -------------------------------------------------

اگر تعداد نمونه‌های سیگنال برابر $N$ و نرخ نمونه‌برداری برابر $f_s$ باشد، مدت‌زمان فایل از رابطه زیر به دست می‌آید:

\begin{equation}
Duration
=
\frac{N}{f_s}
\end{equation}

مدت‌زمان می‌تواند اطلاعات تکمیلی درباره ساختار نمونه‌های گفتاری فراهم کند؛ بااین‌حال، تفسیر اثر آن باید با توجه به نحوه گردآوری داده انجام شود.


% =================================================
\section{الگوریتم‌های طبقه‌بندی}
% =================================================

الگوریتم‌های طبقه‌بندی با توجه به نوع نمایش داده در دو مسیر انتخاب شده‌اند. در مسیر متن، خانواده‌های
\lr{K-Nearest Neighbors}،
\lr{Decision Tree}،
\lr{Logistic Regression}
و
\lr{Multinomial Naive Bayes}
بررسی شده‌اند. در مسیر گفتار نیز خانواده‌های
\lr{Logistic Regression}،
\lr{Support Vector Machine}،
\lr{Random Forest}
و
\lr{K-Nearest Neighbors}
استفاده شده‌اند. در این فصل، مبانی نظری این روش‌ها بیان می‌شود و جزئیات تنظیمات آن‌ها در فصل چهارم ارائه خواهد شد.


% -------------------------------------------------
\subsection{\lr{Logistic Regression}}
% -------------------------------------------------

رگرسیون لجستیک\footnote{\lr{Logistic Regression}} با وجود نام آن، یک الگوریتم طبقه‌بندی است. در حالت چندکلاسه، احتمال تعلق نمونه $x$ به کلاس $k$ را می‌توان با تابع
\lr{Softmax}
نمایش داد:

\begin{equation}
P(y=k\mid x)
=
\frac{
e^{w_k^Tx+b_k}
}{
\sum_{j=1}^{K}e^{w_j^Tx+b_j}
}
\end{equation}

کلاس دارای بیشترین احتمال به‌عنوان خروجی انتخاب می‌شود:

\begin{equation}
\hat{y}
=
\arg\max_k P(y=k\mid x)
\end{equation}

برای کنترل پیچیدگی مدل می‌توان از منظم‌سازی استفاده کرد. شدت منظم‌سازی بر تعادل میان برازش داده‌های آموزشی و توانایی تعمیم مدل اثر می‌گذارد.


% -------------------------------------------------
\subsection{\lr{Multinomial Naive Bayes}}
% -------------------------------------------------

روش بیز ساده یک الگوریتم احتمالاتی مبتنی بر قضیه بیز است:

\begin{equation}
P(C\mid X)
=
\frac{
P(X\mid C)P(C)
}{
P(X)
}
\end{equation}

در این روش فرض می‌شود ویژگی‌ها با شرط مشخص‌بودن کلاس، مستقل از یکدیگر هستند. اگرچه این فرض در بسیاری از داده‌های واقعی دقیق نیست، مدل می‌تواند در طبقه‌بندی متن عملکرد مناسبی داشته باشد.

مدل
\lr{Multinomial Naive Bayes}
برای ویژگی‌های نامنفی مبتنی بر فراوانی یا وزن‌های متنی مناسب است. پارامتر هموارسازی $\alpha$ از صفرشدن احتمال ویژگی‌هایی جلوگیری می‌کند که در نمونه‌های آموزشی یک کلاس مشاهده نشده‌اند.


% -------------------------------------------------
\subsection{\lr{Decision Tree}}
% -------------------------------------------------

درخت تصمیم\footnote{\lr{Decision Tree}} با تقسیم پیاپی فضای ویژگی‌ها، نمونه‌ها را به گره‌هایی با خلوص بیشتر هدایت می‌کند. یکی از معیارهای متداول برای ارزیابی ناخالصی یک گره، شاخص جینی است:

\begin{equation}
Gini
=
1-
\sum_{k=1}^{K}p_k^2
\end{equation}

که در آن $p_k$ نسبت نمونه‌های کلاس $k$ در گره است. درخت‌های بسیار عمیق می‌توانند داده آموزشی را بیش‌ازحد یاد بگیرند؛ بنابراین، کنترل پیچیدگی درخت برای کاهش بیش‌برازش اهمیت دارد.


% -------------------------------------------------
\subsection{\lr{Support Vector Machine}}
% -------------------------------------------------

ماشین بردار پشتیبان یا
\lr{SVM}\footnote{\lr{Support Vector Machine}}
الگوریتمی است که تلاش می‌کند مرز تصمیمی با بیشترین حاشیه میان کلاس‌ها پیدا کند. در حالت خطی، تابع تصمیم به‌صورت زیر تعریف می‌شود:

\begin{equation}
f(x)
=
w^Tx+b
\end{equation}

برای داده‌هایی که کاملاً جداشدنی نیستند، از حاشیه نرم یا
\lr{Soft Margin}
استفاده می‌شود. تابع هدف آن را می‌توان به‌صورت زیر نوشت:

\begin{equation}
\min_{w,b,\xi}
\left(
\frac{1}{2}\left\|w\right\|_2^2
+
C\sum_{i=1}^{N}\xi_i
\right)
\end{equation}

با محدودیت‌های

\begin{equation}
y_i\left(w^T\phi(x_i)+b\right)
\geq
1-\xi_i,
\qquad
\xi_i\geq0
\end{equation}

تعریف می‌شود. پارامتر $C$ تعادل میان بزرگ‌بودن حاشیه و جریمه خطاهای آموزشی را کنترل می‌کند. مقدار کوچک‌تر $C$ منظم‌سازی بیشتری ایجاد می‌کند و مقدار بزرگ‌تر خطاهای آموزشی را با شدت بیشتری جریمه می‌کند.

برای ایجاد مرز غیرخطی می‌توان از تابع کرنل استفاده کرد. کرنل
\lr{RBF}\footnote{\lr{Radial Basis Function}}
به‌صورت زیر تعریف می‌شود:

\begin{equation}
K(x_i,x_j)
=
\exp
\left(
-\gamma
\left\|x_i-x_j\right\|_2^2
\right)
\end{equation}

پارامتر $\gamma$ محدوده اثر هر نمونه را کنترل می‌کند. مقدار کوچک‌تر آن معمولاً مرز هموارتری ایجاد می‌کند، درحالی‌که مقدار بزرگ‌تر می‌تواند به مرزی محلی‌تر و پیچیده‌تر منجر شود.


% -------------------------------------------------
\subsection{\lr{K-Nearest Neighbors}}
% -------------------------------------------------

الگوریتم
\lr{KNN}\footnote{\lr{K-Nearest Neighbors}}
کلاس یک نمونه جدید را براساس برچسب نزدیک‌ترین نمونه‌های آموزشی تعیین می‌کند. نتیجه این مدل به تعداد همسایه‌ها، معیار فاصله و روش وزن‌دهی آرای همسایه‌ها وابسته است.

فاصله
\lr{Minkowski}
به‌صورت زیر تعریف می‌شود:

\begin{equation}
d(x_i,x_j)
=
\left(
\sum_{m=1}^{p}
|x_{im}-x_{jm}|^q
\right)^{1/q}
\end{equation}

که در آن $p$ تعداد ویژگی‌ها و $q$ مرتبه فاصله است. با $q=1$ فاصله منهتن و با $q=2$ فاصله اقلیدسی به دست می‌آید.

برای داده‌های متنی، فاصله کسینوسی نیز کاربرد دارد:

\begin{equation}
d_{\mathrm{cosine}}(x_i,x_j)
=
1-
\frac{x_i^Tx_j}
{\left\|x_i\right\|_2\left\|x_j\right\|_2}
\end{equation}

این معیار بیشتر به جهت بردارها توجه می‌کند و برای نمایش‌های متنی پربعد مناسب است. انتخاب معیار فاصله باید با نوع و مقیاس نمایش ویژگی‌ها سازگار باشد.


% -------------------------------------------------
\subsection{\lr{Random Forest}}
% -------------------------------------------------

جنگل تصادفی\footnote{\lr{Random Forest}} یک روش تجمیعی\footnote{\lr{Ensemble}} مبتنی بر مجموعه‌ای از درخت‌های تصمیم است. هر درخت با بخشی از نمونه‌ها و زیرمجموعه‌ای از ویژگی‌ها ساخته می‌شود و خروجی نهایی طبقه‌بندی از رأی درخت‌ها به دست می‌آید:

\begin{equation}
\hat{y}
=
\mathrm{mode}
\{h_1(x),h_2(x),\ldots,h_M(x)\}
\end{equation}

تجمیع چند درخت می‌تواند واریانس مدل و حساسیت آن به نویز را نسبت به یک درخت منفرد کاهش دهد. تعداد درخت‌ها، عمق آن‌ها و حداقل تعداد نمونه‌های هر برگ از عوامل مؤثر بر پیچیدگی مدل هستند.


% =================================================
\section{الگوریتم‌های خوشه‌بندی}
% =================================================

در هر دو مسیر پروژه، خانواده‌های
\lr{K-Means}،
\lr{Agglomerative Clustering}،
\lr{Gaussian Mixture Model}
و
\lr{DBSCAN}
بررسی شده‌اند. برچسب واقعی زبان نباید در ساخت نمایش بدون نظارت، کاهش ابعاد، انتخاب پارامترها یا برازش الگوریتم خوشه‌بندی استفاده شود. برچسب‌ها فقط پس از تشکیل خوشه‌ها و برای ارزیابی خارجی و تفسیر نتایج قابل استفاده‌اند.


% -------------------------------------------------
\subsection{\lr{K-Means}}
% -------------------------------------------------

الگوریتم
\lr{K-Means}
داده‌ها را به $K$ خوشه تقسیم می‌کند و مجموع مربعات فاصله نمونه‌ها از مرکز خوشه مربوط را کاهش می‌دهد:

\begin{equation}
J
=
\sum_{k=1}^{K}
\sum_{x_i\in C_k}
\left\|x_i-\mu_k\right\|_2^2
\end{equation}

که در آن $C_k$ خوشه شماره $k$ و $\mu_k$ مرکز آن است. این روش به مقداردهی اولیه و تعداد خوشه‌ها حساس است و بیشتر برای خوشه‌هایی با ساختار نسبتاً فشرده و کروی مناسب است.


% -------------------------------------------------
\subsection{\lr{Agglomerative Clustering}}
% -------------------------------------------------

در خوشه‌بندی سلسله‌مراتبی تجمعی، ابتدا هر نمونه یک خوشه مستقل است. سپس در هر مرحله، دو خوشه براساس یک معیار اتصال با یکدیگر ادغام می‌شوند تا ساختاری سلسله‌مراتبی ایجاد شود.

در روش اتصال
\lr{Ward}،
ادغامی انتخاب می‌شود که کمترین افزایش را در مجموع مربعات انحرافات درون خوشه‌ای ایجاد کند. روش‌های اتصال دیگر می‌توانند فاصله میان اعضا یا مراکز خوشه‌ها را به شکل متفاوتی تعریف کنند. انتخاب روش اتصال بر شکل خوشه‌های نهایی اثر می‌گذارد.


% -------------------------------------------------
\subsection{\lr{Gaussian Mixture Model}}
% -------------------------------------------------

مدل مخلوط گاوسی فرض می‌کند توزیع داده از ترکیب چند توزیع گاوسی تشکیل شده است:

\begin{equation}
p(x)
=
\sum_{k=1}^{K}
\pi_k
\mathcal{N}(x\mid\mu_k,\Sigma_k)
\end{equation}

که در آن:

\begin{itemize}

	\item $\pi_k$ وزن مؤلفه $k$؛

	\item $\mu_k$ بردار میانگین مؤلفه؛

	\item $\Sigma_k$ ماتریس کوواریانس آن است.

\end{itemize}

برخلاف
\lr{K-Means}،
این مدل می‌تواند احتمال تعلق هر نمونه به مؤلفه‌های مختلف را محاسبه کند. برای تبدیل این احتمال‌ها به تخصیص قطعی خوشه، مؤلفه دارای بیشترین احتمال انتخاب می‌شود.

یکی از معیارهای انتخاب پیچیدگی مدل مخلوط گاوسی، معیار اطلاعات بیزی یا
\lr{BIC}\footnote{\lr{Bayesian Information Criterion}}
است:

\begin{equation}
BIC
=
-2\log(\hat{\mathcal{L}})
+
q\log(N)
\end{equation}

که در آن $\hat{\mathcal{L}}$ بیشینه درست‌نمایی، $q$ تعداد پارامترهای مدل و $N$ تعداد نمونه‌ها است. مقدار کمتر
\lr{BIC}
نشان‌دهنده تعادل مناسب‌تر میان برازش و پیچیدگی مدل است.


% -------------------------------------------------
\subsection{\lr{DBSCAN}}
% -------------------------------------------------

الگوریتم
\lr{DBSCAN}\footnote{\lr{Density-Based Spatial Clustering of Applications with Noise}}
یک روش خوشه‌بندی مبتنی بر چگالی است. دو پارامتر اصلی آن شعاع همسایگی $\varepsilon$ و حداقل تعداد نقاط همسایه یا
\lr{MinPts}
هستند.

همسایگی نقطه $x$ به‌صورت زیر تعریف می‌شود:

\begin{equation}
N_{\varepsilon}(x)
=
\{x_j\mid d(x,x_j)\leq\varepsilon\}
\end{equation}

یک نقطه زمانی نقطه مرکزی محسوب می‌شود که:

\begin{equation}
|N_{\varepsilon}(x)|
\geq
MinPts
\end{equation}

این الگوریتم می‌تواند خوشه‌هایی با شکل‌های نامنظم را شناسایی کند و نمونه‌های واقع در نواحی کم‌چگالی را به‌عنوان نویز مشخص سازد. عملکرد آن به انتخاب مقیاس ویژگی‌ها و پارامترهای چگالی حساس است.


% =================================================
\section{آماده‌سازی و کاهش ابعاد ویژگی‌ها}
% =================================================

پیش از اجرای بعضی الگوریتم‌ها، لازم است نمایش ویژگی‌ها از نظر مقیاس، مقدارهای گمشده و ابعاد آماده شود. این مراحل باید در ارزیابی نظارت‌شده فقط براساس داده آموزشی برازش شوند تا اطلاعات بخش ارزیابی وارد فرایند یادگیری نشود.


% -------------------------------------------------
\subsection{جایگزینی مقادیر گمشده}
% -------------------------------------------------

یکی از راه‌های مدیریت مقادیر گمشده در ویژگی‌های عددی، جایگزینی آن‌ها با یک آماره مقاوم مانند میانه ویژگی است:

\begin{equation}
x_{ij}
\leftarrow
\mathrm{median}(X_{\mathrm{train},j})
\end{equation}

میانه باید فقط از داده آموزشی محاسبه شود؛ زیرا محاسبه آن با استفاده از کل داده می‌تواند اطلاعات بخش ارزیابی را وارد مرحله پیش‌پردازش کند.


% -------------------------------------------------
\subsection{استانداردسازی ویژگی‌ها}
% -------------------------------------------------

ویژگی‌های عددی ممکن است مقیاس‌های متفاوتی داشته باشند. این تفاوت به‌ویژه بر الگوریتم‌های مبتنی بر فاصله و بعضی مرزهای تصمیم اثر می‌گذارد. استانداردسازی یک ویژگی به‌صورت زیر تعریف می‌شود:

\begin{equation}
z
=
\frac{x-\mu}{\sigma}
\end{equation}

که در آن $\mu$ میانگین و $\sigma$ انحراف معیار ویژگی در داده آموزشی است. پس از تبدیل، ویژگی‌ها در داده آموزشی دارای میانگین نزدیک به صفر و انحراف معیار نزدیک به یک خواهند بود. مدل‌های مبتنی بر درخت معمولاً نسبت به مقیاس ویژگی‌ها حساسیت کمتری دارند.


% -------------------------------------------------
\subsection{\lr{Principal Component Analysis}}
% -------------------------------------------------

تحلیل مؤلفه‌های اصلی یا
\lr{PCA}\footnote{\lr{Principal Component Analysis}}
یک روش خطی برای تبدیل ویژگی‌های هم‌بسته به مؤلفه‌هایی متعامد است. اگر ماتریس کوواریانس داده با $\Sigma$ نمایش داده شود، بردارها و مقادیر ویژه از رابطه زیر به دست می‌آیند:

\begin{equation}
\Sigma v_i
=
\lambda_i v_i
\end{equation}

که در آن $v_i$ بردار ویژه و $\lambda_i$ مقدار ویژه متناظر است. مؤلفه‌هایی با مقدار ویژه بزرگ‌تر، سهم بیشتری از واریانس داده را توضیح می‌دهند. با نگه‌داشتن تعدادی از مؤلفه‌های نخست می‌توان نمایش فشرده‌تری از داده‌های عددی ساخت.

در این پروژه، این روش برای ایجاد نمایش بدون نظارت ویژگی‌های گفتاری در مسیر خوشه‌بندی به کار می‌رود. تعداد مؤلفه‌ها و نحوه اجرای آن در فصل سوم تشریح خواهند شد.


% -------------------------------------------------
\subsection{\lr{Truncated SVD}}
% -------------------------------------------------

ماتریس‌های متنی
\lr{TF--IDF}
معمولاً پربعد و تنک هستند. تجزیه مقدار منفرد به‌صورت زیر تعریف می‌شود:

\begin{equation}
X
=
U\Sigma V^T
\end{equation}

در نسخه بریده‌شده، تنها تعداد محدودی از مؤلفه‌ها نگه داشته می‌شوند:

\begin{equation}
X
\approx
U_r\Sigma_rV_r^T
\end{equation}

روش
\lr{Truncated SVD}
بدون نیاز به مرکزدهی کامل ماتریس تنک، یک نمایش کم‌بعد از داده ایجاد می‌کند. در این پروژه، این روش در ساخت نمایش بدون نظارت متن برای خوشه‌بندی استفاده می‌شود. تعداد دقیق مؤلفه‌ها در فصل سوم ارائه خواهد شد.


% -------------------------------------------------
\subsection{نرمال‌سازی بردارها}
% -------------------------------------------------

در نرمال‌سازی
$L_2$،
هر بردار بر اندازه اقلیدسی خود تقسیم می‌شود:

\begin{equation}
\hat{x}
=
\frac{x}{\left\|x\right\|_2}
\end{equation}

پس از این تبدیل، طول هر بردار برابر یک می‌شود و جهت بردار نقش بیشتری در مقایسه نمونه‌ها پیدا می‌کند. این تبدیل برای نمایش کم‌بعد متن در مسیر خوشه‌بندی مناسب است.


% =================================================
\section{تنظیم و اعتبارسنجی مدل}
% =================================================

ارزیابی معتبر یک مدل نیازمند جداسازی داده ارزیابی نهایی از فرایند آموزش و تنظیم ابرپارامترها است. همچنین، اگر چند نمونه به یک گوینده یا منبع مشترک تعلق داشته باشند، قرارگرفتن آن‌ها در دو بخش آموزش و ارزیابی می‌تواند به برآورد خوش‌بینانه عملکرد منجر شود.


% -------------------------------------------------
\subsection{تقسیم گروه‌آگاه داده}
% -------------------------------------------------

در تقسیم گروه‌آگاه، تمام نمونه‌های متعلق به یک گروه مشترک در یک بخش قرار می‌گیرند. اگر مجموعه گروه‌های آموزش و آزمون به‌ترتیب با $G_{\mathrm{train}}$ و $G_{\mathrm{test}}$ نمایش داده شوند، شرط جلوگیری از هم‌پوشانی عبارت است از:

\begin{equation}
G_{\mathrm{train}}
\cap
G_{\mathrm{test}}
=
\varnothing
\end{equation}

در حالت طبقه‌بندی، بهتر است علاوه بر استقلال گروه‌ها، توزیع کلاس‌ها نیز تا حد امکان در بخش‌های مختلف حفظ شود. تعریف گروه‌ها و اندازه بخش‌های داده در فصل سوم ارائه خواهند شد.


% -------------------------------------------------
\subsection{\lr{Cross Validation}}
% -------------------------------------------------

در اعتبارسنجی متقابل، داده آموزشی به چند بخش تقسیم می‌شود و مدل چند بار آموزش می‌بیند. در هر تکرار، یک بخش برای اعتبارسنجی و سایر بخش‌ها برای آموزش استفاده می‌شوند.

اگر امتیاز مدل در هر تکرار برابر $S_i$ باشد، میانگین امتیاز اعتبارسنجی از رابطه زیر به دست می‌آید:

\begin{equation}
CV
=
\frac{1}{K}
\sum_{i=1}^{K}S_i
\end{equation}

در اعتبارسنجی گروه‌آگاه، نمونه‌های یک گروه میان بخش آموزش و اعتبارسنجی تقسیم نمی‌شوند. این موضوع احتمال نشت اطلاعات ناشی از وابستگی نمونه‌ها را کاهش می‌دهد. تعداد بخش‌های اعتبارسنجی در فصل چهارم مشخص خواهد شد.


% -------------------------------------------------
\subsection{\lr{Grid Search}}
% -------------------------------------------------

در جست‌وجوی شبکه‌ای، مجموعه‌ای از مقادیر ممکن برای هر ابرپارامتر تعریف می‌شود. تمام ترکیب‌های حاصل با استفاده از اعتبارسنجی متقابل ارزیابی می‌شوند و ترکیبی که معیار انتخاب را بهینه کند، انتخاب می‌شود.

داده آزمون نهایی نباید برای انتخاب ابرپارامتر یا خانواده مدل استفاده شود. پس از پایان تنظیم، مدل منتخب روی داده آموزشی برازش شده و فقط برای ارزیابی نهایی روی داده آزمون اجرا می‌شود.


% -------------------------------------------------
\subsection{جلوگیری از نشت اطلاعات}
% -------------------------------------------------

نشت اطلاعات زمانی رخ می‌دهد که اطلاعاتی از داده اعتبارسنجی یا آزمون، مستقیم یا غیرمستقیم وارد فرایند آموزش شود. نمونه‌هایی از نشت اطلاعات عبارت‌اند از:

\begin{itemize}

	\item برازش استانداردسازی یا جایگزینی مقادیر گمشده روی کل داده؛

	\item ساخت واژگان متن با استفاده از اسناد آزمون؛

	\item انتخاب ویژگی با استفاده از برچسب‌های داده آزمون؛

	\item حضور نمونه‌های یک گوینده یا منبع مشترک در آموزش و آزمون؛

	\item استفاده از برچسب زبان در برازش یا انتخاب پارامترهای خوشه‌بندی.

\end{itemize}

قرار دادن مراحل پیش‌پردازش نظارت‌شده در زنجیره آموزش و حفظ استقلال گروه‌ها، از مهم‌ترین روش‌های کنترل این خطر هستند.


% =================================================
\section{معیارهای ارزیابی طبقه‌بندی}
% =================================================

در مسئله چندکلاسه، معیارها ابتدا برای هر کلاس و سپس با روش‌های مختلف میانگین‌گیری محاسبه می‌شوند. برای یک کلاس مشخص، $TP$ تعداد مثبت‌های درست، $FP$ تعداد مثبت‌های نادرست و $FN$ تعداد منفی‌های نادرست است.


% -------------------------------------------------
\subsection{\lr{Accuracy}}
% -------------------------------------------------

دقت کلی نسبت تعداد پیش‌بینی‌های صحیح به کل نمونه‌ها است:

\begin{equation}
Accuracy
=
\frac{1}{N}
\sum_{i=1}^{N}
\mathbb{I}(y_i=\hat{y}_i)
\end{equation}

اگر $C$ ماتریس درهم‌ریختگی باشد، فرم چندکلاسه آن عبارت است از:

\begin{equation}
Accuracy
=
\frac{
\sum_{k=1}^{K}C_{kk}
}{
\sum_{i=1}^{K}\sum_{j=1}^{K}C_{ij}
}
\end{equation}


% -------------------------------------------------
\subsection{\lr{Precision}}
% -------------------------------------------------

معیار
\lr{Precision}
نشان می‌دهد از میان نمونه‌هایی که مدل به یک کلاس نسبت داده است، چه نسبتی واقعاً متعلق به آن کلاس بوده‌اند:

\begin{equation}
Precision
=
\frac{TP}{TP+FP}
\end{equation}


% -------------------------------------------------
\subsection{\lr{Recall}}
% -------------------------------------------------

معیار
\lr{Recall}
نشان می‌دهد چه نسبتی از نمونه‌های واقعی یک کلاس به‌درستی تشخیص داده شده‌اند:

\begin{equation}
Recall
=
\frac{TP}{TP+FN}
\end{equation}


% -------------------------------------------------
\subsection{\lr{F1-Score}}
% -------------------------------------------------

معیار
\lr{F1-Score}
میانگین هارمونیک
\lr{Precision}
و
\lr{Recall}
است:

\begin{equation}
F1
=
2
\frac{
Precision\times Recall
}{
Precision+Recall
}
\end{equation}


% -------------------------------------------------
\subsection{میانگین‌های \lr{Macro} و \lr{Weighted}}
% -------------------------------------------------

در میانگین
\lr{Macro}،
تمام کلاس‌ها وزن یکسانی دارند:

\begin{equation}
MacroPrecision
=
\frac{1}{K}
\sum_{k=1}^{K}Precision_k
\end{equation}

\begin{equation}
MacroRecall
=
\frac{1}{K}
\sum_{k=1}^{K}Recall_k
\end{equation}

\begin{equation}
MacroF1
=
\frac{1}{K}
\sum_{k=1}^{K}F1_k
\end{equation}

در میانگین وزنی، سهم هر کلاس متناسب با تعداد نمونه‌های واقعی آن است. برای مثال:

\begin{equation}
WeightedF1
=
\sum_{k=1}^{K}
\frac{n_k}{N}F1_k
\end{equation}

که در آن $n_k$ تعداد نمونه‌های کلاس $k$ و $N$ تعداد کل نمونه‌ها است. معیارهای
\lr{Macro}
برای بررسی عملکرد متوازن روی زبان‌ها و معیارهای وزنی برای نمایش عملکرد متناسب با فراوانی کلاس‌ها مفید هستند.


% -------------------------------------------------
\subsection{\lr{Confusion Matrix}}
% -------------------------------------------------

ماتریس درهم‌ریختگی\footnote{\lr{Confusion Matrix}} زبان واقعی نمونه‌ها را با زبان پیش‌بینی‌شده مقایسه می‌کند. اگر $C_{ij}$ عضو سطر $i$ و ستون $j$ باشد، داریم:

\begin{equation}
C_{ij}
=
\text{تعداد نمونه‌های کلاس واقعی }i
\text{ که به کلاس }j
\text{ نسبت داده شده‌اند}
\end{equation}

مقادیر قطر اصلی نشان‌دهنده پیش‌بینی‌های صحیح و مقادیر خارج از قطر بیانگر خطاهای میان کلاس‌ها هستند.


% =================================================
\section{معیارهای ارزیابی خوشه‌بندی}
% =================================================

معیارهای خوشه‌بندی به دو دسته اصلی تقسیم می‌شوند:

\begin{itemize}

	\item معیارهای داخلی که فقط از ویژگی‌ها و تخصیص خوشه‌ها استفاده می‌کنند؛

	\item معیارهای خارجی که تخصیص خوشه‌ها را پس از برازش با برچسب‌های واقعی مقایسه می‌کنند.

\end{itemize}

معیارهای داخلی می‌توانند برای انتخاب بدون نظارت به کار روند. معیارهای خارجی در این پروژه فقط برای ارزیابی پس از برازش و تفسیر نتایج استفاده می‌شوند.


% -------------------------------------------------
\subsection{\lr{Silhouette Score}}
% -------------------------------------------------

معیار
\lr{Silhouette Score}
انسجام درون خوشه‌ای و جدایی میان خوشه‌ها را هم‌زمان بررسی می‌کند. برای نمونه $i$، اگر $a(i)$ میانگین فاصله آن با اعضای خوشه خودش و $b(i)$ کمترین میانگین فاصله آن با یک خوشه دیگر باشد، داریم:

\begin{equation}
s(i)
=
\frac{
b(i)-a(i)
}{
\max\{a(i),b(i)\}
}
\end{equation}

و

\begin{equation}
-1\leq s(i)\leq1
\end{equation}

مقدار نزدیک به یک نشان‌دهنده تخصیص مناسب‌تر، مقدار نزدیک به صفر نشان‌دهنده قرارگرفتن در مرز خوشه‌ها و مقدار منفی نشان‌دهنده تخصیص نامناسب احتمالی است.


% -------------------------------------------------
\subsection{\lr{Calinski--Harabasz Index}}
% -------------------------------------------------

معیار
\lr{Calinski--Harabasz}
نسبت پراکندگی بین خوشه‌ها به پراکندگی درون خوشه‌ها را اندازه‌گیری می‌کند:

\begin{equation}
CH
=
\frac{
\mathrm{tr}(B_K)/(K-1)
}{
\mathrm{tr}(W_K)/(N-K)
}
\end{equation}

که در آن $B_K$ ماتریس پراکندگی بین خوشه‌ای، $W_K$ ماتریس پراکندگی درون خوشه‌ای، $K$ تعداد خوشه‌ها و $N$ تعداد نمونه‌ها است. مقدار بزرگ‌تر معمولاً نشان‌دهنده خوشه‌های فشرده‌تر و جداشده‌تر است.


% -------------------------------------------------
\subsection{\lr{Davies--Bouldin Index}}
% -------------------------------------------------

معیار
\lr{Davies--Bouldin}
شباهت هر خوشه با شبیه‌ترین خوشه دیگر را بررسی می‌کند:

\begin{equation}
DB
=
\frac{1}{K}
\sum_{i=1}^{K}
\max_{j\neq i}
\left(
\frac{S_i+S_j}{M_{ij}}
\right)
\end{equation}

که در آن $S_i$ پراکندگی درون خوشه $i$ و $M_{ij}$ فاصله میان مراکز خوشه‌های $i$ و $j$ است. مقدار کمتر این معیار نشان‌دهنده کیفیت بهتر خوشه‌بندی است.


% -------------------------------------------------
\subsection{\lr{Cluster Purity}}
% -------------------------------------------------

معیار
\lr{Purity}
بررسی می‌کند اعضای هر خوشه تا چه اندازه به یک کلاس واقعی تعلق دارند. اگر مجموعه خوشه‌ها برابر

\begin{equation}
\Omega=\{\omega_1,\ldots,\omega_K\}
\end{equation}

و مجموعه کلاس‌های واقعی برابر $C=\{c_1,\ldots,c_J\}$ باشد، داریم:

\begin{equation}
Purity(\Omega,C)
=
\frac{1}{N}
\sum_k
\max_j
|\omega_k\cap c_j|
\end{equation}

مقدار بیشتر نشان‌دهنده همگنی زبانی بیشتر خوشه‌ها است. بااین‌حال، این معیار به افزایش تعداد خوشه‌ها حساس است و می‌تواند با ایجاد خوشه‌های کوچک‌تر افزایش یابد؛ بنابراین، بهتر است در کنار معیارهای دیگر تفسیر شود.


% -------------------------------------------------
\subsection{\lr{Adjusted Rand Index}}
% -------------------------------------------------

معیار
\lr{Adjusted Rand Index}
توافق میان تخصیص خوشه‌ها و برچسب‌های واقعی را با اصلاح توافق مورد انتظار در حالت تصادفی اندازه‌گیری می‌کند:

\begin{equation}
ARI
=
\frac{
RI-E[RI]
}{
\max(RI)-E[RI]
}
\end{equation}

مقدار یک نشان‌دهنده انطباق کامل است. مقدار نزدیک به صفر بیانگر توافقی مشابه حالت تصادفی است و مقدار منفی می‌تواند نشان‌دهنده توافق کمتر از انتظار تصادفی باشد.


% -------------------------------------------------
\subsection{\lr{Normalized Mutual Information}}
% -------------------------------------------------

معیار
\lr{NMI}\footnote{\lr{Normalized Mutual Information}}
میزان اطلاعات مشترک میان تخصیص خوشه‌ها و برچسب‌های واقعی را اندازه‌گیری می‌کند. در نرمال‌سازی حسابی می‌توان نوشت:

\begin{equation}
NMI(C,Y)
=
\frac{
2I(C;Y)
}{
H(C)+H(Y)
}
\end{equation}

که در آن $I(C;Y)$ اطلاعات متقابل و $H(C)$ و $H(Y)$ آنتروپی خوشه‌ها و برچسب‌های واقعی هستند. مقدار این معیار در بازه صفر تا یک قرار می‌گیرد و مقدار بزرگ‌تر نشان‌دهنده انطباق بیشتر است.


% =================================================
\section{خلاصه و جمع‌بندی}
% =================================================

در این فصل، مبانی نظری مورد نیاز برای شناسایی زبان از داده‌های متنی و گفتاری بررسی شد. ابتدا تفاوت یادگیری نظارت‌شده و بدون نظارت و نقش آن‌ها در طبقه‌بندی و خوشه‌بندی معرفی شد.

در مسیر متن، مفاهیم
\lr{N-Gram}،
\lr{TF--IDF}
و انتخاب ویژگی با
\lr{Chi-Square}
توضیح داده شدند. در مسیر گفتار نیز ویژگی‌های
\lr{MFCC}،
مشتق‌های زمانی، نرخ عبور از صفر، ویژگی‌های طیفی، انرژی و مدت‌زمان معرفی شدند.

سپس مبانی الگوریتم‌های طبقه‌بندی و خوشه‌بندی مورد استفاده در پروژه بررسی شدند. مفاهیم مدیریت مقادیر گمشده، استانداردسازی، کاهش ابعاد، نرمال‌سازی، اعتبارسنجی متقابل، تنظیم ابرپارامترها و جلوگیری از نشت اطلاعات نیز تشریح شدند.

در پایان، معیارهای ارزیابی طبقه‌بندی و خوشه‌بندی معرفی شدند. در خوشه‌بندی، معیارهای داخلی بدون استفاده از برچسب واقعی محاسبه می‌شوند و معیارهای خارجی فقط پس از تشکیل خوشه‌ها برای تفسیر نتایج به کار می‌روند.

جزئیات مجموعه‌داده‌ها، مراحل پیش‌پردازش و ساختار نهایی ویژگی‌ها در فصل سوم و نحوه پیاده‌سازی، تنظیم مدل‌ها و نتایج آزمایش‌ها در فصل چهارم ارائه خواهند شد.
